\documentclass[12pt]{article}

\usepackage[spanish]{babel}
\usepackage[utf8]{inputenc}
\usepackage[T1]{fontenc}
\usepackage{geometry}
\usepackage{graphicx}
\usepackage{float}
\usepackage{booktabs}
\usepackage{longtable}
\usepackage{array}
\usepackage{enumitem}
\usepackage{hyperref}
\usepackage{titlesec}
\usepackage{setspace}
\usepackage{caption}

\geometry{
    letterpaper,
    margin=2.5cm
}

\setstretch{1.2}

\title{
    \textbf{Entregable 2: Avance de Agente e Investigación} \\
    \large PF-3311 Agentes Virtuales Inteligentes
}


\hyphenpenalty=10000
\exhyphenpenalty=10000
\sloppy


\author{
    Luis Carlos Quesada Rodríguez
}

\date{I Ciclo 2026}

\begin{document}

\maketitle

\tableofcontents
\newpage

% =========================================================
\section{Resumen de Cambios desde el Entregable 1}

\subsection{Feedback Recibido}

Durante la evaluación del Entregable 1 se indicó que las preguntas RQ1 y RQ2 se encontraban correctamente formuladas y abordaban constructos distintos, destacando especialmente el enfoque de detección implícita de familiaridad conductual. También se señaló que la decisión de mantener el mismo avatar entre condiciones para aislar el efecto del comportamiento era metodológicamente adecuada.

Sin embargo, se indicó que las preguntas RQ3 y RQ4 carecían de instrumentos estandarizados claramente asociados para su medición. Como recomendación, se sugirió utilizar el instrumento SAM (Self-Assessment Manikin) para medir respuestas afectivas y utilizar ítems relacionados con \textit{perceived knowledge} en agentes virtuales para fortalecer la RQ4.

\subsection{Cambios Realizados}

A partir de la retroalimentación recibida, se realizaron ajustes principalmente en la operacionalización de las preguntas de investigación y en el diseño metodológico del estudio.

Las preguntas RQ3 y RQ4 fueron refinadas para alinearlas de manera más clara con constructos medibles y con instrumentos respaldados por literatura previa en HCI y agentes conversacionales. Además, se incorporó explícitamente el uso del instrumento SAM para medir respuesta afectiva y escalas relacionadas con \textit{perceived knowledge} para evaluar percepción de conocimiento contextual previo.

También se refinó el enfoque conceptual del proyecto. La propuesta evolucionó desde una idea más cercana a imitación directa hacia un enfoque centrado en \textit{familiaridad conductual}, entendida como la reproducción parcial de patrones conversacionales reconocibles sin intentar replicar completamente la identidad de una persona.

En términos técnicos, la arquitectura evolucionó hacia un sistema basado en LangGraph, incorporando perfiles conductuales, memoria contextual, mecanismos de refinamiento iterativo y validación piloto de perfiles antes de la ejecución experimental.

\subsection{Justificación de los Cambios}

Los cambios realizados tuvieron como principal objetivo fortalecer el anclaje metodológico y teórico de la propuesta.

La reformulación parcial de las preguntas de investigación permitió establecer relaciones más claras entre constructos, métricas e instrumentos de medición. Asimismo, el cambio hacia un enfoque de familiaridad conductual permitió reducir posibles problemas relacionados con antropomorfismo excesivo o dependencia de similitud visual.

Finalmente, la incorporación de mecanismos de validación y refinamiento busca asegurar que los perfiles conductuales generados produzcan señales suficientemente consistentes de familiaridad antes de realizar la evaluación experimental principal.







% =========================================================
\section{Estado Actual del Agente}

\subsection{Descripción General del Prototipo}

El estado actual del proyecto corresponde a una prueba de concepto funcional orientada a la evaluación de familiaridad conductual en agentes virtuales inteligentes. El sistema implementa un entorno conversacional basado en Godot, donde un participante interactúa con un agente virtual mediante texto y síntesis de voz.

Actualmente, el sistema permite ejecutar sesiones experimentales utilizando dos condiciones diferenciadas:

\begin{itemize}
\item una condición basada en perfiles conductuales personalizados;
\item una condición de control basada en un perfil conversacional genérico.
\end{itemize}

Ambas condiciones utilizan el mismo avatar, voz e interfaz visual con el objetivo de aislar el efecto de los patrones conductuales sobre la percepción del usuario.

La arquitectura del sistema utiliza un backend en Python basado en FastAPI y LangGraph para la orquestación conversacional, generación de respuestas y administración de perfiles conductuales. La generación conversacional se encuentra integrada actualmente mediante modelos locales ejecutados con Ollama, mientras que la síntesis de voz utiliza Edge-TTS.

Además, el sistema incorpora mecanismos de entrenamiento, validación y refinamiento de perfiles conductuales orientados a modelar patrones conversacionales de una persona conocida sin intentar replicar explícitamente su identidad.

\subsection{Arquitectura Actual del Sistema}

La arquitectura actual del sistema sigue un modelo cliente-servidor desacoplado orientado a la ejecución de sesiones experimentales controladas. El cliente, desarrollado en Godot, concentra la experiencia interactiva del participante, incluyendo la interfaz conversacional, la representación visual del agente mediante un avatar VRM, la reproducción de voz sintetizada y la navegación del protocolo experimental.

Por otra parte, el backend implementado en Python utiliza FastAPI y LangGraph para coordinar la lógica conversacional y experimental del sistema. Esta capa es responsable de la generación de respuestas mediante modelos de lenguaje, la administración de perfiles conductuales, la síntesis de voz y el almacenamiento de logs experimentales.

Actualmente, el sistema incorpora múltiples grafos especializados dentro de LangGraph, orientados a distintas tareas del flujo experimental. Entre ellos se incluyen grafos para interacción conversacional, entrenamiento de perfiles conductuales, refinamiento iterativo y validación piloto de perfiles antes de la ejecución experimental principal.

La comunicación entre Godot y el backend se realiza mediante WebSockets para la interacción conversacional en tiempo real y mediante HTTP para operaciones relacionadas con configuración experimental, entrenamiento y validación de perfiles.

\begin{table}[H]
\centering
\small
\renewcommand{\arraystretch}{1.2}

\begin{tabular}{p{4cm} p{9cm}}
\toprule
\textbf{Componente} & \textbf{Responsabilidad Principal} \\
\midrule

Godot
& Interfaz del participante, representación visual del agente, navegación experimental y reproducción de audio. \\

FastAPI
& Comunicación cliente-servidor, manejo de sesiones y exposición de rutas experimentales. \\

LangGraph
& Orquestación de flujos conversacionales, entrenamiento y validación de perfiles conductuales. \\

Ollama
& Generación conversacional mediante modelos de lenguaje locales. \\

Edge-TTS
& Síntesis de voz para respuestas del agente. \\

Sistema de perfiles conductuales
& Construcción y aplicación de patrones conversacionales personalizados. \\

SQLite y almacenamiento local
& Persistencia de sesiones, logs y perfiles experimentales. \\

\bottomrule
\end{tabular}

\caption{Resumen de componentes principales del sistema.}

\end{table}

\begin{figure}[H]
    \centering
    % \includegraphics[width=0.9\textwidth]{architecture.png}
    \caption{Arquitectura general actual del sistema.}
\end{figure}

\subsection{Estado de Implementación}

La Tabla 2 resume el estado actual de las principales funcionalidades implementadas dentro del prototipo, distinguiendo entre componentes completamente funcionales, parcialmente implementados y elementos aún pendientes dentro del flujo experimental.

\begin{table}[H]
\centering
\small
\renewcommand{\arraystretch}{1.2}

\begin{tabular}{p{6cm} p{2.5cm} p{5cm}}
\toprule
\textbf{Funcionalidad} & \textbf{Estado} & \textbf{Observaciones} \\
\midrule

Interfaz conversacional en Godot
& Funcional
& Flujo experimental y chat implementados. \\

Sesiones experimentales A/B
& Funcional
& Contrabalanceo intra-sujetos operativo. \\

Integración LangGraph
& Funcional
& Grafos para chat, entrenamiento y validación. \\

Generación conversacional
& Funcional
& Integración mediante Ollama. \\

Síntesis de voz (Edge-TTS)
& Funcional
& Reproducción de audio integrada en Godot. \\

Avatar VRM y animaciones
& Funcional
& Lip-sync y reacciones básicas implementadas. \\

Sistema de perfiles conductuales
& Funcional
& Entrenamiento y perfiles YAML operativos. \\

Refinamiento iterativo de perfiles
& Parcial
& Flujo backend implementado; integración UI pendiente. \\

Validación piloto de perfiles
& Funcional
& Evaluación y métricas iniciales disponibles. \\

Entrada por voz (Whisper/STT)
& Opcional
& Integración futura aún no prioritaria para el estudio. \\

Cuestionarios integrados
& Pendiente
& Actualmente externos al sistema. \\

\bottomrule
\end{tabular}

\caption{Estado actual de implementación del prototipo.}

\end{table}

\subsection{Retos Actuales de Implementación}

Actualmente, el sistema ya permite generar perfiles conductuales mediante entrevistas conversacionales guiadas por modelos de lenguaje, así como procesos de refinamiento y validación iterativa de dichos perfiles. Experiencias previas realizadas fuera del entorno experimental actual han mostrado resultados satisfactorios en la construcción de perfiles capaces de reproducir patrones conversacionales reconocibles y consistentes.

No obstante, uno de los principales objetivos actuales del proyecto consiste en integrar estos mecanismos dentro de la experiencia interactiva desarrollada en Godot, manteniendo simultáneamente criterios de naturalidad conversacional, usabilidad y control experimental. En particular, el proyecto busca adaptar y mejorar técnicas previamente utilizadas en otros entornos para que puedan funcionar de manera más orgánica dentro del flujo experimental propuesto.

Otro reto importante corresponde al balance entre flexibilidad conversacional y consistencia metodológica. El sistema busca permitir conversaciones suficientemente naturales para capturar estilos narrativos, mecanismos de seguimiento conversacional y patrones de respuesta del usuario modelado, manteniendo al mismo tiempo un nivel de control adecuado para la comparación experimental entre condiciones.

Finalmente, algunos componentes complementarios del sistema continúan en evaluación o refinamiento, incluyendo integración opcional de entrada por voz y mecanismos más robustos de validación de perfiles conductuales antes de la ejecución de estudios piloto.


\subsection{Capturas y Diagramas}

Para la versión actual del prototipo se consideran relevantes las siguientes capturas y diagramas:

\begin{itemize}
\item menú experimental principal;
\item interfaz conversacional del agente;
\item selección de orden experimental A/B;
\item entrenamiento de perfiles conductuales;
\item validación piloto de perfiles;
\item arquitectura general del sistema;
\item ejemplo de perfil conductual estructurado.
\end{itemize}


% =========================================================
\section{Diseño Metodológico del Estudio}

\subsection{Tipo de Estudio}

El presente trabajo corresponde a una evaluación piloto orientada a validar la viabilidad técnica y metodológica de una prueba de concepto basada en familiaridad conductual dentro de agentes virtuales inteligentes.

Debido a las restricciones éticas y metodológicas asociadas a evaluaciones formales con usuarios reales, el estudio se plantea principalmente como una instancia exploratoria y de refinamiento experimental, enfocada en analizar la percepción de familiaridad conductual dentro de interacciones conversacionales breves.

La evaluación busca validar tanto el funcionamiento técnico del prototipo como la consistencia de las señales conductuales generadas por el agente antes de una posible evaluación formal posterior.



\subsection{Participantes}

La evaluación piloto utilizará una muestra pequeña de participantes reclutados mediante conveniencia. Debido a la naturaleza del estudio, los participantes corresponderán principalmente a personas cercanas al entorno social del investigador y de la persona utilizada para la construcción del perfil conductual, incluyendo amistades, conocidos y contactos de fácil acceso para la ejecución de las sesiones piloto.

La selección de participantes busca favorecer que los usuarios posean suficiente familiaridad previa con la persona modelada como para potencialmente reconocer patrones conductuales durante la interacción, sin revelar explícitamente la existencia de perfiles de familiaridad dentro del experimento.

Se espera contar con entre 6 y 12 participantes para las sesiones piloto iniciales, permitiendo obtener retroalimentación cualitativa y métricas preliminares relacionadas con percepción de familiaridad, naturalidad y respuesta afectiva.

Adicionalmente, el protocolo incorpora una etapa previa de validación de perfiles conductuales. En esta fase, un subconjunto independiente de participantes evaluará la consistencia y recognoscibilidad de los perfiles generados antes de la ejecución principal del experimento. Estos participantes actuarán únicamente como evaluadores piloto del perfil conductual y no participarán posteriormente en las sesiones experimentales principales, con el objetivo de evitar sesgos derivados de exposición previa al sistema.

Como criterios de inclusión se consideran:

\begin{itemize}
    \item ser mayor de edad;
    \item poseer experiencia básica utilizando sistemas conversacionales digitales;
    \item mantener algún nivel de familiaridad previa con la persona utilizada para construir el perfil conductual;
    \item comprender el idioma español utilizado durante la interacción.
\end{itemize}

Como criterios de exclusión se consideran:

\begin{itemize}
    \item dificultades importantes de comunicación durante la interacción;
    \item conocimiento técnico directo sobre la implementación específica de las condiciones experimentales;
    \item conocimiento explícito previo sobre los objetivos específicos relacionados con familiaridad conductual dentro del estudio;
    \item haber participado previamente en la etapa de validación piloto de perfiles conductuales.
\end{itemize}

Con el objetivo de reducir posibles sesgos de expectativa o priming, los participantes no serán informados explícitamente sobre la naturaleza exacta de la manipulación experimental ni sobre la existencia de perfiles conductuales modelados durante la interacción. Las instrucciones experimentales serán presentadas de forma neutral, enfocándose únicamente en la interacción conversacional con el agente virtual.

El reclutamiento se realizará mediante invitación directa dentro del entorno social y académico cercano al proyecto.







\subsection{Diseño Experimental}

El estudio utilizará un diseño intra-sujetos con contrabalanceo parcial orientado a evaluar el efecto de patrones de familiaridad conductual sobre la percepción del usuario durante interacciones conversacionales breves.

El protocolo experimental estará compuesto por dos etapas diferenciadas:

\begin{itemize}
    \item una etapa previa de validación piloto de perfiles conductuales;
    \item una etapa principal de interacción experimental entre participantes y el agente virtual.
\end{itemize}


\subsubsection{Validación Piloto de Perfiles Conductuales}

Antes de la ejecución principal del experimento, los perfiles conductuales generados serán sometidos a una etapa piloto de validación orientada a verificar la consistencia y recognoscibilidad de los patrones conversacionales modelados.

Esta etapa se fundamenta parcialmente en trabajos relacionados con percepción de familiaridad y \textit{perceived knowledge} en agentes virtuales, particularmente los discutidos por Poletaykina et al. (2025). Sin embargo, en contraste con enfoques centrados únicamente en percepción subjetiva de familiaridad, esta propuesta incorpora además una fase explícita de validación metodológica orientada a verificar que los perfiles conductuales produzcan señales suficientemente consistentes antes de su utilización dentro de la evaluación principal.

Durante esta fase, un subconjunto independiente de participantes familiarizados con la persona utilizada para construir el perfil conductual interactuará con el agente virtual en sesiones breves de conversación utilizando la Condición A. Posteriormente, los evaluadores completarán un conjunto de escalas Likert de 1 a 7 orientadas específicamente a evaluar calidad y recognoscibilidad del perfil generado.

Las dimensiones evaluadas incluyen:

\begin{itemize}
    \item similitud de tono conversacional;
    \item similitud de expresiones y redacción;
    \item similitud en longitud y estructura de respuestas;
    \item consistencia conductual;
    \item percepción general de familiaridad;
    \item naturalidad conversacional;
    \item ausencia de filtración explícita de identidad.
\end{itemize}

A partir de estas evaluaciones, el sistema calcula métricas agregadas relacionadas con similitud conversacional, naturalidad e integridad de identidad. Los perfiles únicamente serán considerados aptos para la evaluación principal cuando alcancen umbrales mínimos predefinidos de consistencia y recognoscibilidad.

Adicionalmente, el protocolo permitirá recopilar retroalimentación cualitativa y correcciones manuales por parte de los evaluadores, incluyendo observaciones relacionadas con expresiones atípicas, inconsistencias narrativas o patrones conversacionales poco representativos de la persona modelada.

Los perfiles que no alcancen niveles suficientes de consistencia o recognoscibilidad serán refinados iterativamente antes de ser utilizados dentro de la evaluación experimental principal.



\subsubsection{Evaluación Experimental Principal}

En la etapa principal, cada participante interactuará con ambas condiciones experimentales:

\begin{itemize}
    \item una condición basada en familiaridad conductual;
    \item una condición conversacional genérica de control.
\end{itemize}

Para reducir posibles efectos de orden, los participantes serán divididos en dos grupos. El primer grupo interactuará inicialmente con la Condición A y posteriormente con la Condición B, mientras que el segundo grupo seguirá el orden inverso.

Cada interacción tendrá una duración aproximada de cinco minutos, seguida de cuestionarios y evaluación posterior relacionados con percepción de familiaridad, naturalidad, presencia social y respuesta afectiva.

\subsection{Condiciones Experimentales}

\subsubsection{Condición Experimental: Condición A}

La condición experimental corresponde a un Agente Virtual Inteligente implementado en Godot, utilizando un avatar antropomórfico genérico y una voz sintética estándar. La familiaridad del agente emerge principalmente a partir de patrones conversacionales asociados a una persona conocida por el participante.

Estos patrones incluyen elementos como tono comunicativo, expresiones recurrentes, estructura de respuesta, mecanismos de seguimiento conversacional y estilo general de interacción. Adicionalmente, el sistema puede incorporar referencias contextuales derivadas del perfil conductual modelado.

Los perfiles utilizados en esta condición serán previamente sometidos a una etapa piloto de validación y refinamiento antes de su incorporación dentro del flujo experimental principal.

\subsubsection{Condición de Control: Condición B}

La condición de control utiliza el mismo sistema, avatar, voz e interfaz conversacional. Sin embargo, el agente emplea comportamiento conversacional genérico y no incorpora patrones de familiaridad conductual ni referencias contextuales asociadas al perfil modelado.

De esta manera, la principal diferencia entre condiciones corresponde específicamente al comportamiento conversacional del agente y al uso de patrones de familiaridad conductual durante la interacción.

\subsubsection{Justificación del Diseño}

La utilización del mismo embodiment, interfaz y entorno conversacional en ambas condiciones busca aislar específicamente el efecto de los patrones conductuales sobre la percepción del usuario.

Asimismo, la incorporación de una etapa previa de validación piloto permite reducir el riesgo de utilizar perfiles conductuales insuficientemente consistentes o reconocibles dentro de la evaluación principal.

De esta manera, el estudio busca analizar si señales de familiaridad, naturalidad, presencia social y conocimiento contextual percibido pueden emerger principalmente a partir del comportamiento conversacional del agente y no de elementos visuales o técnicos externos.




\bibliographystyle{IEEEtran}
\bibliography{references}

\end{document}
